\documentclass[10pt,a4paper]{amsart}
\usepackage{amsmath,amssymb,amsthm,graphicx}
% Retained minimal preamble for the Stacks Project chapter more-morphisms.tex (More on Morphisms).
% Verbatim excerpt of the chapter preamble: Xy-pic setup (native preamble lines 14-18) and the native
% theorem declarations with the mathematical macros (native preamble lines 154-217). Copied mechanically;
% nothing added, removed or redefined. The wrapper class and the amsmath/amssymb/amsthm/graphicx base
% packages are supplied by the assembler, so the class/externaldocument part of the native preamble is omitted.
% For commutative diagrams we use Xy-pic
\usepackage[all]{xy}

% We use 2cell for 2-commutative diagrams.
\xyoption{2cell}
% Theorem environments.
%
\theoremstyle{plain}
\newtheorem{theorem}[subsection]{Theorem}
\newtheorem{proposition}[subsection]{Proposition}
\newtheorem{lemma}[subsection]{Lemma}

\theoremstyle{definition}
\newtheorem{definition}[subsection]{Definition}
\newtheorem{example}[subsection]{Example}
\newtheorem{exercise}[subsection]{Exercise}
\newtheorem{situation}[subsection]{Situation}

\theoremstyle{remark}
\newtheorem{remark}[subsection]{Remark}
\newtheorem{remarks}[subsection]{Remarks}

\numberwithin{equation}{subsection}

% Macros
%
\def\lim{\mathop{\mathrm{lim}}\nolimits}
\def\colim{\mathop{\mathrm{colim}}\nolimits}
\def\Spec{\mathop{\mathrm{Spec}}}
\def\Hom{\mathop{\mathrm{Hom}}\nolimits}
\def\Ext{\mathop{\mathrm{Ext}}\nolimits}
\def\SheafHom{\mathop{\mathcal{H}\!\mathit{om}}\nolimits}
\def\SheafExt{\mathop{\mathcal{E}\!\mathit{xt}}\nolimits}
\def\Sch{\mathit{Sch}}
\def\Mor{\mathop{\mathrm{Mor}}\nolimits}
\def\Ob{\mathop{\mathrm{Ob}}\nolimits}
\def\Sh{\mathop{\mathit{Sh}}\nolimits}
\def\NL{\mathop{N\!L}\nolimits}
\def\CH{\mathop{\mathrm{CH}}\nolimits}
\def\proetale{{pro\text{-}\acute{e}tale}}
\def\etale{{\acute{e}tale}}
\def\QCoh{\mathit{QCoh}}
\def\Ker{\mathop{\mathrm{Ker}}}
\def\Im{\mathop{\mathrm{Im}}}
\def\Coker{\mathop{\mathrm{Coker}}}
\def\Coim{\mathop{\mathrm{Coim}}}

% Boxtimes
%
\DeclareMathSymbol{\boxtimes}{\mathbin}{AMSa}{"02}

%
% Macros for moduli stacks/spaces
%
\def\QCohstack{\mathcal{QC}\!\mathit{oh}}
\def\Cohstack{\mathcal{C}\!\mathit{oh}}
\def\Spacesstack{\mathcal{S}\!\mathit{paces}}
\def\Quotfunctor{\mathrm{Quot}}
\def\Hilbfunctor{\mathrm{Hilb}}
\def\Curvesstack{\mathcal{C}\!\mathit{urves}}
\def\Polarizedstack{\mathcal{P}\!\mathit{olarized}}
\def\Complexesstack{\mathcal{C}\!\mathit{omplexes}}
% \Pic is the operator that assigns to X its picard group, usage \Pic(X)
% \Picardstack_{X/B} denotes the Picard stack of X over B
% \Picardfunctor_{X/B} denotes the Picard functor of X over B
\def\Pic{\mathop{\mathrm{Pic}}\nolimits}
\def\Picardstack{\mathcal{P}\!\mathit{ic}}
\def\Picardfunctor{\mathrm{Pic}}
\def\Deformationcategory{\mathcal{D}\!\mathit{ef}}
\newenvironment{reference}{}{}
\newenvironment{slogan}{}{}
\newenvironment{history}{}{}
\usepackage{cleveref}
\crefname{originallemma}{original lemma}{original lemmas}
\Crefname{originallemma}{Original lemma}{Original lemmas}
\makeatletter
\newlabel{lemma-slice-smooth}{{057G}{source}{}{sourcepointer-1}{}}
\newlabel{lemma-slice-smooth@cref}{{[originallemma][0][]057G}{[1][0][]source}{}{}{}}
\makeatother
\crefname{originallemma}{original lemma}{original lemmas}
\Crefname{originallemma}{Original lemma}{Original lemmas}
\makeatletter
\newlabel{varieties-lemma-smooth-separable-closed-points-dense}{{056U}{source}{}{sourcepointer-2}{}}
\newlabel{varieties-lemma-smooth-separable-closed-points-dense@cref}{{[originallemma][0][]056U}{[1][0][]source}{}{}{}}
\makeatother
\title{Every point of the image of a smooth morphism admits an etale local section}
\author{The Stacks Project Authors}
\date{stacks-project:more-morphisms.tex:local-snapshot-20261011; GNU FDL 1.2 or later}
\begin{document}
\maketitle
\noindent\textit{Editorial note.} Statement span is lines 10845-10857 of sources/stacks-more-morphisms/pinned-original-source.tex; proof span is lines 10859-10868, the complete author proof (8 body lines).\par\medskip
\noindent\textit{Editorial note.} Navigation stubs: 2 cited-but-not-included labels are supplied as non-load-bearing pointers with their pinned tags where the snapshot tag index has them.\par\medskip
\noindent\textit{Editorial note.} extra\_packages is empty on purpose: the retained preamble already loads xy with options and a second declaration clashes in the assembler wrapper.\par\medskip
\noindent\textit{Editorial note.} Text is retained verbatim, including the author's own typos and placeholders.\par\medskip
\noindent\textit{Editorial note.} Only the author's FIRST proof (lines 10859-10868) is retained; the chapter contains a second proof immediately below, outside the collected spans.\par\medskip
\noindent\textit{Editorial note.} The statement's reference environment cites \textbackslash{}cite[Corollaire 17.16.3 (ii)]\{EGA4\}; a bibliography.tex holding the snapshot my.bib EGA4 entry verbatim is attached (precedent S3-morph-bertini-irreducible, CONTROL section 8.1).\par\medskip
\noindent\textit{Allowed standard prerequisites (editorial).}\begin{itemize}
\item Explicit prerequisite boundary. The lemma that smooth fibres over a field have dense closed points with finite separable residue fields (Varieties) and the slice lemma producing an etale locally closed slice through such a point (same chapter); cited inputs with the two-step argument written in the retained span.
\end{itemize}\medskip
\ifdefined\hypertarget\hypertarget{sourcepointer-1}{}\fi
\noindent\textit{Source reference.} Original lemma 057G: sources/stacks-more-morphisms/pinned-original-source.tex, source lines 10775--10783. Cited by the retained statement or proof but not part of the collected spans; supplied only as a navigation pointer. The argument does not depend on its content.\par\medskip
\ifdefined\hypertarget\hypertarget{sourcepointer-2}{}\fi
\noindent\textit{Source reference.} Original lemma 056U: sources/stacks-more-morphisms/pinned-original-source.tex, source lines 10868--10868. Cited by the retained statement or proof but not part of the collected spans; supplied only as a navigation pointer. The argument does not depend on its content.\par\medskip
\begin{lemma}
\label{lemma-etale-nbhd-dominates-smooth}
\begin{slogan}
Smooth morphisms admit \'etale local sections.
\end{slogan}
\begin{reference}
See \cite[Corollaire 17.16.3 (ii)]{EGA4}.
\end{reference}
Let $f : X \to S$ be a smooth morphism of schemes.
Let $s \in S$ be a point in the image of $f$.
Then there exists an \'etale neighbourhood $(S', s') \to (S, s)$
and a $S$-morphism $S' \to X$.
\end{lemma}

\begin{proof}[First proof of Lemma \ref{lemma-etale-nbhd-dominates-smooth}]
By assumption $X_s \not = \emptyset$. By
Varieties, Lemma \ref{varieties-lemma-smooth-separable-closed-points-dense}
there exists a closed point $x \in X_s$ such that $\kappa(x)$
is a finite separable field extension of $\kappa(s)$.
Hence by
Lemma \ref{lemma-slice-smooth}
there exists an immersion $Z \to X$ such that $Z \to S$ is \'etale and such
that $x \in Z$. Take $(S' , s') = (Z, x)$.
\end{proof}

\begin{thebibliography}{1}
\bibitem{EGA4}
Grothendieck, Alexander and Dieudonn{\'e}, Jean.
\newblock {\em {{\'E}l{\'e}ments de g{\'e}om{\'e}trie alg{\'e}brique {IV}}}.
\newblock Publications {M}ath{\'e}matiques 20, 24, 28, 32. Institute des {H}autes {\'E}tudes {S}cientifiques, 1964--1967.
\end{thebibliography}
\end{document}
